\documentclass[10pt,twocolumn]{article}

\usepackage[margin=0.70in]{geometry}
\usepackage{amsmath,amssymb,mathtools}
\usepackage{microtype}
\usepackage{times}
\usepackage[hidelinks]{hyperref}
\usepackage{enumitem}

\setlist{nosep,leftmargin=*}
\setlength{\columnsep}{0.22in}
\setlength{\parindent}{1em}
\setlength{\parskip}{0pt}
\emergencystretch=1em

\title{\vspace{-1.2em}Takeoff Kernels for Recursive Self-Improvement:\\
A Short Formal Sketch}
\author{Jonathan R. Landers}
\date{July 2026}

\begin{document}
\maketitle
\vspace{-1.5em}

\begin{abstract}
Current theories of language-model self-improvement usually ask whether a
generator can use its own verification signal to produce a better successor.
That question records the amount or eventual limit of improvement, but not the
finite shape by which improvement appears across rounds.  This sketch applies
the takeoff-kernel viewpoint to the standard generator--verifier--update loop.
For bounded self-training, the kernel records what fraction of the eventual
gain arrives in each round.  For sustained recursive improvement, it records
how log-capability growth approaches its terminal velocity.  The resulting
language distinguishes immediate, delayed, monotone, ringing, mixed, and
diffusive takeoff even when two systems have the same terminal performance or
asymptotic growth rate.
\end{abstract}

\section{Context}

Recent self-improvement theory studies a recurring pattern.  A model generates
candidate answers, a verifier scores or filters those answers, and the accepted
data are distilled back into a new model.  The verifier may be external, or it
may be another use of the same model.  Sharpening theory interprets the process
as moving probability mass toward outputs the model can recognize as better
\cite{huang2025}.  The generation--verification gap measures whether such
reweighting has positive value \cite{song2025}.  More recent analyses study the
multi-round trajectory, finite-sample effects, saturation, and easy-to-hard
curricula \cite{sun2026,liu2026}.

This is usually bounded self-training, not unrestricted recursive
self-improvement.  Nevertheless, it already defines a discrete response curve
over improvement rounds.  The takeoff-kernel question is: not only \emph{does}
performance rise, but \emph{when and in what shape} does the rise occur?

\section{Standard self-improvement loop}

Let $\mathcal Q$ be a question or task space, $P$ a task distribution, and
$\mathcal A$ an answer space.  At round $t$:
\begin{align}
q&\sim P, & a&\sim\pi_t(\cdot\mid q),
\end{align}
where $\pi_t$ is the current solver.  Let
$s_t(q,a)\in[0,1]$ be a verifier score and let $w$ be a nonnegative filtering
or weighting function.  The verifier-induced target distribution is
\begin{equation}
\bar\pi_t(a\mid q)
=\frac{\pi_t(a\mid q)w(s_t(q,a))}{Z_t(q)},
\end{equation}
where $Z_t(q)$ normalizes the distribution.  Supervised fine-tuning or
reinforcement learning then implements an approximate update
\begin{equation}
\pi_{t+1}=\mathcal U(\pi_t,\bar\pi_t;B_t),
\end{equation}
with sample and compute budget $B_t$.

Let $r(q,a)$ denote the evaluation utility one ultimately cares about, which
need not equal the verifier score, and define
\begin{equation}
J(\pi)=\mathbb E_{q\sim P,\,a\sim\pi(\cdot\mid q)}[r(q,a)],
\qquad J_t=J(\pi_t).
\end{equation}
The idealized generation--verification advantage is
\begin{equation}
g_t=J(\bar\pi_t)-J(\pi_t).
\end{equation}
A positive $g_t$ says that verification exposes better probability mass.  It
does not say how efficiently finite training transfers that advantage into the
successor.  Schematically,
\begin{equation}
J_{t+1}-J_t\approx \eta_t g_t-\ell_t,
\end{equation}
where $\eta_t$ is uptake efficiency and $\ell_t$ collects sampling,
optimization, distribution-shift, and verifier-mismatch losses.

\section{Bounded improvement kernel}

Suppose the repeated update converges to $J_\infty>J_0$.  Define the normalized
takeoff profile
\begin{equation}
F_t^{\mathrm b}=\frac{J_t-J_0}{J_\infty-J_0},
\qquad F_0^{\mathrm b}=0,
\end{equation}
and its causal finite difference
\begin{equation}
\kappa_t^{\mathrm b}=F_t^{\mathrm b}-F_{t-1}^{\mathrm b},
\qquad t\ge1.
\end{equation}
Then
\begin{equation}
\sum_{t\ge1}\kappa_t^{\mathrm b}=1.
\end{equation}
Thus $\kappa_t^{\mathrm b}$ is the fraction of eventual capability gain
delivered in round $t$.  If performance improves monotonically, the kernel is
nonnegative and can be read as a probability distribution over improvement
time.  If some entries are negative, the response is signed: the system
overshoots, regresses, or rings before settling.

Two useful shape statistics transfer directly.  For tolerance $\varepsilon$,
the settling time is
\begin{equation}
T_\varepsilon=\min\{T:|F_t^{\mathrm b}-1|\le\varepsilon
\text{ for all }t\ge T\},
\end{equation}
and the excess signed response is
\begin{equation}
R(\kappa^{\mathrm b})=\sum_{t\ge1}|\kappa_t^{\mathrm b}|-1.
\end{equation}
The first measures how long improvement remains in startup; the second measures
instability invisible in the final score.

\section{Sustained recursive takeoff}

Bounded benchmark accuracy cannot itself sustain a positive exponential
velocity.  For open-ended recursive improvement, choose an unbounded effective
capability coordinate $C_t>0$: for example, the size of the reliably solvable
task frontier, inverse error $(1-J_t)^{-1}$, or success odds where meaningful.
Define the capability velocity
\begin{equation}
u_t=\log C_{t+1}-\log C_t.
\end{equation}
If $u_t\to\alpha>0$, define
\begin{equation}
F_t^{\mathrm r}=\frac{u_t}{\alpha},
\qquad
\kappa_t^{\mathrm r}=F_t^{\mathrm r}-F_{t-1}^{\mathrm r}.
\end{equation}
This is the direct analogue of a velocity takeoff kernel for recursive
redundancy \cite{landers2026}.  The terminal number $\alpha$ says how fast
capability eventually grows; the kernel says how that growth rate becomes
visible.  Therefore two systems can share the same $\alpha$ while having
arbitrarily different finite takeoff.  A genuinely accelerating regime in
which $u_t$ itself grows without bound would require a second-order extension
rather than normalization by a fixed terminal velocity.

\section{Takeoff types and interpretation}

The same vocabulary now classifies self-improvement trajectories.
\begin{itemize}
\item \emph{No takeoff.} Verification supplies no usable advantage, or learning
cannot absorb it.
\item \emph{Immediate.} Most kernel mass lies in the first rounds; a strong verifier
quickly extracts the available gain.
\item \emph{Delayed.} Early rounds contribute little until coverage, curriculum, or
verification quality crosses a threshold.
\item \emph{Monotone long tail.} Gains stay positive but decay as the solver catches
its verifier or the supply of accepted information is exhausted.
\item \emph{Ringing.} Negative kernel entries mark alternating improvement and
regression from reward bias, distribution shift, or evaluator drift.
\item \emph{Mixed.} A fast output-level channel is followed by slower weight, tool,
architecture, or meta-agent improvement.
\item \emph{Diffusive.} Algebraic rather than exponential settling suggests a broad
or growing state space rather than finitely many dominant modes.
\end{itemize}

The generation--verification gap is therefore best viewed as an available
forcing signal; the takeoff kernel is the realized temporal response.  Under a
local linear approximation, sequential stages such as verification, filtering,
and distillation compose by convolution, while parallel feedback channels mix.
Outside that approximation the update is nonlinear, but the measured kernel
still gives a model-independent summary of onset, delay, settling, and
instability.

The basic separation principle is simple: terminal gain $J_\infty-J_0$, or
terminal log-capability velocity $\alpha$, does not determine takeoff.  A fuller
evaluation of recursive self-improvement should report the response shape as
well: its kernel, settling time, signed response, and, where a modal fit is
appropriate, its slowest settling mode.

\begin{thebibliography}{9}\footnotesize

\bibitem{huang2025}
A.~Huang, A.~Block, D.~J. Foster, et al.
\newblock Self-improvement in language models: The sharpening mechanism.
\newblock \emph{ICLR}, 2025. \href{https://arxiv.org/abs/2412.01951}{arXiv:2412.01951}.

\bibitem{song2025}
Y.~Song, H.~Zhang, C.~Eisenach, S.~Kakade, D.~Foster, and U.~Ghai.
\newblock Mind the gap: Examining the self-improvement capabilities of large language models.
\newblock \emph{ICLR}, 2025. \href{https://arxiv.org/abs/2412.02674}{arXiv:2412.02674}.

\bibitem{sun2026}
Y.~Sun, Y.~Liang, Z.~Zhang, X.~Liu, and J.~Teng.
\newblock Theoretical modeling of large language model self-improvement training dynamics through solver--verifier gap.
\newblock \emph{ICLR}, 2026. \href{https://arxiv.org/abs/2507.00075}{arXiv:2507.00075}.

\bibitem{liu2026}
C.~Liu, Y.~Dong, Y.~Shen, and Q.~Lei.
\newblock A task-centric theory for iterative self-improvement with easy-to-hard curricula.
\newblock \emph{ICLR}, 2026. \href{https://arxiv.org/abs/2602.10014}{arXiv:2602.10014}.

\bibitem{landers2026}
J.~R. Landers.
\newblock The shape of recursive redundancy: Velocity takeoff kernels for memoization.
\newblock Manuscript, 2026.

\end{thebibliography}

\end{document}
